% Source PDF: source/普通高考/2010/2010北京文.pdf, page 1 (question 9).
% The source chart is vector artwork (no embedded bitmap); node text and topology
% are transcribed from the rendered page.
% Node-edge list: 开始 -> 输入 x -> ①; ①-(是)-> y=2-x -> 输出 y;
% ①-(否)-> ② -> merge before 输出 y; 输出 y -> 结束.
% The two branches merge at the single connector immediately above 输出 y;
% no other edges or crossings are present in the source.
\begin{tikzpicture}[
  baseline,
  >=stealth,
  x=1cm,
  y=1.25cm,
  line width=0.55pt,
  line cap=round,
  line join=round,
  font=\normalsize,
  every node/.style={
    anchor=center,
    align=center,
    execute at begin node={\setlength{\parindent}{0pt}},
  },
  flowbox/.style={
    draw,
    minimum height=0.66cm,
    inner xsep=4pt,
    inner ysep=2pt,
    align=center,
    execute at begin node={\setlength{\parindent}{0pt}},
  },
  terminal/.style={flowbox,rounded corners=0.18cm,minimum width=1.25cm},
  io/.style={
    flowbox,
    shape=trapezium,
    trapezium left angle=75,
    trapezium right angle=105,
    minimum width=2.35cm,
    minimum height=0.72cm,
  },
  decision/.style={
    flowbox,
    shape=diamond,
    aspect=2.8,
    inner sep=0pt,
    minimum width=3.65cm,
    minimum height=1.00cm,
  },
]
  % Main vertical path.
  \node[terminal] (start) at (0,8.95) {开始};
  \node[io] (input) at (0,7.65) {输入 $x$};
  \node[decision] (check) at (0,6.20) {\textcircled{1}};
  \node[flowbox,minimum width=2.35cm] (lower) at (0,4.78) {$y=2-x$};
  \node[io] (output) at (0,3.35) {输出 $y$};
  \node[terminal] (finish) at (0,1.95) {结束};

  % False branch to the right-hand assignment box.
  \node[flowbox,minimum width=2.35cm] (upper) at (3.12,4.78) {\textcircled{2}};

  \draw[->] (start.south) -- (input.north);
  \draw[->] (input.south) -- (check.north);
  \draw[->] (check.south) -- node[left=2pt] {是} (lower.north);
  \draw[->] (lower.south) -- (output.north);
  \draw[->] (check.east) -- (3.12,6.20) -- (upper.north);
  \draw[->] (upper.south) -- (3.12,3.90) -- (0,3.90);
  \draw[->] (output.south) -- (finish.north);

  \node[anchor=south west,inner sep=1pt] at (1.00,6.36) {否};
  \path[use as bounding box] (-1.95,9.35) rectangle (4.25,1.65);
\end{tikzpicture}
